\editionPage{indices-defis}{Indications progressives}{Lire une aide, puis reprendre la recherche}{Les indices orientent la démarche ; les solutions restent dans la partie Corrigés.}
\begin{blocEdition}{indices-defis}
\editionEtape{D01 · Le dernier souffle de César}
\textbf{Indice 1.} Convertir la masse de l'atmosphère en grammes avant d'utiliser $n=\dfrac{m}{M}$. Pour une inspiration, utiliser
le volume molaire : $n=\dfrac{V}{V_\mathrm{m}}$.\par\smallskip
\textbf{Indice 2.} Le dernier souffle représente la fraction $\dfrac{n_\text{souffle}}{n_\text{atmosphère}}$ de l'atmosphère. Si le mélange est uniforme,
chacune de nos inspirations contient cette même fraction des molécules du dernier souffle.\par\smallskip
\editionEtape{D02 · Deux colorants, une seule couleur}
\textbf{Indice 1.} À \qty{638}{nm}, un seul des deux colorants absorbe : son absorbance donne directement sa concentration.\par\smallskip
\textbf{Indice 2.} Reporter $c_\mathrm{B}$ dans l'expression de $A_{427}$ : on obtient une équation du premier degré en $c_\mathrm{J}$. Pour la
question 6, résoudre par substitution, comme un système de deux équations à deux inconnues.\par\smallskip
\editionEtape{D03 · L'arbre de Diane prend du poids}
\textbf{Indice 1.} Chaque atome de cuivre qui réagit quitte le fil, mais deux atomes d'argent viennent s'y déposer. Pour
trouver le réactif épuisé, comparer $n(\ce{Ag+})$ à $2\,n(\ce{Cu})$.\par\smallskip
\textbf{Indice 2.} Exprimer la masse de cuivre perdue et la masse d'argent gagnée en fonction de $x$, puis les combiner.
Pour la dernière question, penser à la couleur de la solution et au chapitre~2.\par\smallskip
\editionEtape{D04 · La sonde lambda}
\textbf{Indice 1.} Ajuster d'abord C et H, puis O ; si le nombre de molécules de dioxygène n'est pas entier, tout
multiplier par 2. Pour l'air, $n_\text{air}=\dfrac{n(\ce{O2})}{0{,}21}$ et $m_\text{air}=n_\text{air}\times\qty{29,0}{g.mol^{-1}}$.\par\smallskip
\textbf{Indice 2.} Pour $\lambda=0{,}90$, la quantité de dioxygène vaut $0{,}90\times12{,}5$ mol. Comparer ensuite les quotients
$\dfrac{n_0}{a}$ des deux réactifs, avec leurs nombres stœchiométriques 2 et 25.\par\smallskip
\editionEtape{D05 · Que vaut vraiment une eau de Javel ?}
\textbf{Indice 1.} Dans un titrage indirect, on enchaîne deux bilans : chaque ion hypochlorite forme une molécule de diiode,
et chaque molécule de diiode réagit avec deux ions thiosulfate.\par\smallskip
\textbf{Indice 2.} Le « chlore actif » est la masse de dichlore que libérerait un volume donné de solution : une mole d'ions
hypochlorite libère une mole de $\ce{Cl2}$, de masse molaire \qty{71,0}{g.mol^{-1}}. Ne pas oublier la dilution.\par\smallskip
\end{blocEdition}

\editionPage{indices-defis-2}{Indications progressives (suite)}{Lire une aide, puis reprendre la recherche}{Les indices orientent la démarche ; les solutions restent dans la partie Corrigés.}
\begin{blocEdition}{indices-defis-2}
\editionEtape{D06 · La molécule d'eau en géométrie}
\textbf{Indice 1.} Le centre du cube est à mi-chemin de deux sommets opposés : la distance C--H est la moitié de la
grande diagonale du cube, que l'on obtient avec le théorème de Pythagore appliqué deux fois.\par\smallskip
\textbf{Indice 2.} Les deux vecteurs ont même norme : leur somme est portée par la bissectrice de l'angle $\alpha$. Projeter
chacun d'eux sur cette bissectrice.\par\smallskip
\editionEtape{D07 · L'étiquette d'une eau minérale}
\textbf{Indice 1.} Une concentration en masse en \unit{mg.L^{-1}} divisée par une masse molaire en \unit{g.mol^{-1}} donne directement une
concentration en \unit{mmol.L^{-1}}. Les électrons ont une masse négligeable : $M(\ce{SO4^{2-}})=M(\ce{S})+4M(\ce{O})$.\par\smallskip
\textbf{Indice 2.} Un ion $\ce{Ca^{2+}}$ apporte deux charges : il compte double. Pour la question 4, comparer les deux concentrations
les plus grandes ; une mole de gypse libère une mole de chacun de ces ions.\par\smallskip
\editionEtape{D08 · L'odeur du chocolat}
\textbf{Indice 1.} Pour énumérer les isomères, partir des chaînes de cinq atomes de carbone, puis de quatre avec un groupe
méthyle, puis de trois avec deux groupes méthyle ; placer le groupe carbonyle à toutes les positions possibles.\par\smallskip
\textbf{Indice 2.} Les deux bandes moyennes entre 2700 et \qty{2850}{cm^{-1}} sont caractéristiques d'une famille. Pour l'indice C,
repérer, dans chaque candidat ramifié, le carbone voisin du groupe carbonyle.\par\smallskip
\editionEtape{D09 · Deux procédés pour l'ibuprofène}
\textbf{Indice 1.} Si une étape transforme \qty{90}{\percent} de ce qu'elle reçoit, l'étape suivante ne reçoit que ces \qty{90}{\percent}, dont elle
transforme à son tour \qty{90}{\percent}. Pour la question 4, une économie d'atomes de \qty{40}{\percent} signifie que \qty{1}{kg} de produit exige
$\dfrac{1}{0{,}40}$ kg de réactifs.\par\smallskip
\textbf{Indice 2.} Masse de réactifs nécessaire $=\dfrac{\text{masse de produit}}{\text{économie d'atomes}\times\text{rendement global}}$.\par\smallskip
\end{blocEdition}

\editionPage{indices-defis-3}{Indications progressives (suite)}{Lire une aide, puis reprendre la recherche}{Les indices orientent la démarche ; les solutions restent dans la partie Corrigés.}
\begin{blocEdition}{indices-defis-3}
\editionEtape{D10 · Le meilleur carburant ?}
\textbf{Indice 1.} $\mathrm{PC}=\dfrac{|E_\mathrm{r}|}{M}$ en \unit{kJ.g^{-1}}, c'est-à-dire en \unit{MJ.kg^{-1}}. L'énergie par litre vaut le pouvoir calorifique multiplié par la
masse volumique.\par\smallskip
\textbf{Indice 2.} Pour le kérosène, $|E_\mathrm{r}|=\mathrm{PC}\times M$. Une mole de $\ce{C12H26}$ donne douze moles de $\ce{CO2}$. Pour la question 5, se demander
ce qui compte le plus quand la fusée doit soulever ses réservoirs depuis le sol, puis quand elle est déjà en altitude.\par\smallskip
\editionEtape{D11 · Ce qui tient le noyau}
\textbf{Indice 1.} Les deux protons portent la charge $+e$ ; on utilise la même distance $d$ pour les deux forces. Pour la question 2,
$P=m\,g$ donne $m=\dfrac{F}{g}$.\par\smallskip
\textbf{Indice 2.} Pour la question 5, isoler $d$ dans la loi de Coulomb : $d=e\,\sqrt{\dfrac{k}{F}}$. Pour la question 6, écrire le rapport
littéral des deux forces : que devient $d^2$ ?\par\smallskip
\editionEtape{D12 · Les hémisphères de Magdebourg}
\textbf{Indice 1.} Convertir $R$ en mètres ; la force vaut la différence de pression multipliée par $\pi\,R^2$. Pour l'arbre, se demander quelle
force exerce le crochet fixé à l'arbre quand un seul attelage tire.\par\smallskip
\textbf{Indice 2.} Mariotte pour l'air intérieur : $P\,V=P'\,(V+v)$. En Python, utiliser une boucle \code{while P > 0.10} qui met à jour \code{P} et
ajoute 1 au compteur à chaque tour.\par\smallskip
\editionEtape{D13 · Freiner à temps}
\textbf{Indice 1.} Pendant la réaction, la vitesse ne change pas : $d=v\times\Delta t$, avec $v$ en \unit{m.s^{-1}}. Au freinage,
$\sum F=0{,}8\,m\,g$, donc $\Delta v=\dfrac{\sum F}{m}\,\Delta t=0{,}8\,g\,\Delta t$.\par\smallskip
\textbf{Indice 2.} Dans la boucle \code{while v > 0}, avancer la position de \code{v * dt}, puis diminuer la vitesse de \code{0.8 * 9.81 * dt}. Pour la règle des
deux secondes, penser que le véhicule de devant freine lui aussi.\par\smallskip
\editionEtape{D14 · Ne pas faire sauter les plombs}
\textbf{Indice 1.} $I=\dfrac{P}{U}$ pour chaque appareil ; sur un même circuit, les intensités s'additionnent. L'abonnement limite la puissance
totale à \qty{6000}{W}, soit une intensité $\dfrac{6000}{230}$.\par\smallskip
\textbf{Indice 2.} \code{from itertools import combinations} donne tous les groupes de $k$ appareils : faire varier $k$ de 1 à 6 et sommer les
puissances de chaque groupe. Pour la rallonge, convertir la section en \unit{m^2} ($\qty{1}{mm^2}=\qty{1e-6}{m^2}$), puis $P=R\,I^2$.\par\smallskip
\end{blocEdition}

\editionPage{indices-defis-4}{Indications progressives (suite)}{Lire une aide, puis reprendre la recherche}{Les indices orientent la démarche ; les solutions restent dans la partie Corrigés.}
\begin{blocEdition}{indices-defis-4}
\editionEtape{D15 · Rouler moins vite pour consommer moins}
\textbf{Indice 1.} À vitesse constante, $\Delta E_\mathrm{c}=0$ : le travail moteur compense exactement le travail des forces résistantes,
$W_\mathrm{moteur}=(f_\mathrm{air}+f_\mathrm{roulement})\times d$ avec $d=\qty{1,00e5}{m}$.\par\smallskip
\textbf{Indice 2.} Un litre fournit $0{,}30\times\qty{35}{MJ}=\qty{10,5}{MJ}$ de travail : diviser le travail moteur par \qty{10,5e6}{J}. Dans le programme, une
boucle \code{for} sur les trois vitesses, sans oublier de les convertir en \unit{m.s^{-1}}.\par\smallskip
\editionEtape{D16 · Stocker l'électricité dans un lac}
\textbf{Indice 1.} $E_\mathrm{p}=m\,g\,h$ avec $m=\rho\,V$. L'énergie restituée vaut $0{,}90\,E_\mathrm{p}$ ; l'énergie de pompage vaut $\dfrac{E_\mathrm{p}}{0{,}85}$ ; le rendement
global est leur rapport.\par\smallskip
\textbf{Indice 2.} Chaque seconde, un débit $Q$ (en \unit{m^3.s^{-1}}) fournit la puissance $0{,}90\,\rho\,Q\,g\,h$. La durée de fonctionnement est le volume
divisé par le débit.\par\smallskip
\editionEtape{D17 · Localiser l'épicentre d'un séisme}
\textbf{Indice 1.} Les ondes P mettent $\dfrac{d}{6{,}0}$ et les ondes S $\dfrac{d}{3{,}5}$ pour parcourir la même distance $d$ : le retard est la différence des deux
durées. L'épicentre est à la distance $d$ de chaque station : il est sur un cercle centré sur la station.\par\smallskip
\textbf{Indice 2.} Pour chaque point $(x\,;\,y)$, calculer l'écart $\sum(d_\text{calculée}-d_\text{mesurée})^2$ sur les trois stations, et garder le point qui donne le plus petit
écart. Deux boucles \code{for} imbriquées parcourent la carte.\par\smallskip
\editionEtape{D18 · Photographier la Lune}
\textbf{Indice 1.} Avec $\alg{OA}=\qty{-3,844e8}{m}$, le terme $\dfrac{1}{\alg{OA}}$ est négligeable devant $\dfrac{1}{f'}$. Le rayon issu du bord de la Lune qui passe par O n'est
pas dévié : appliquer le théorème de Thalès avec $\alg{OA'}\approx f'$.\par\smallskip
\textbf{Indice 2.} Garder les distances dans la même unité pour calculer $\alpha$ (sans unité). Pour la dernière question, résoudre
$f'\times\alpha=\qty{24}{mm}$.\par\smallskip
\editionEtape{D19 · Pourquoi le ciel est-il bleu ?}
\textbf{Indice 1.} Le rapport des intensités diffusées vaut $\left(\dfrac{\lambda_\text{rouge}}{\lambda_\text{bleu}}\right)^4$. La lumière qui nous parvient du ciel, hors de la direction du Soleil,
est de la lumière diffusée.\par\smallskip
\textbf{Indice 2.} Si l'épaisseur d'atmosphère vaut $e$ à la verticale, un rayon incliné de \ang{5} au-dessus de l'horizon en traverse environ $\dfrac{e}{\cos\ang{85}}$. Plus le
trajet est long, plus le bleu est retiré de la lumière directe.\par\smallskip
\end{blocEdition}

\editionPage{indices-defis-5}{Indications progressives (suite)}{Lire une aide, puis reprendre la recherche}{Les indices orientent la démarche ; les solutions restent dans la partie Corrigés.}
\begin{blocEdition}{indices-defis-5}
\editionEtape{D20 · De quoi les étoiles sont-elles faites ?}
\textbf{Indice 1.} Une raie sombre est une radiation absorbée : calculer $E_5-E_2$, puis $\lambda=\dfrac{h\,c}{E_5-E_2}$ en convertissant les eV en joules.
Le programme peut parcourir les raies mesurées, puis, pour chacune, tous les éléments et toutes leurs raies (trois boucles imbriquées).\par\smallskip
\textbf{Indice 2.} Les raies visibles de l'hydrogène partent toutes du niveau $n=2$ : il faut donc que des atomes s'y trouvent déjà. Pour la
raie de 1868, comparer l'écart avec la raie la plus proche au critère de \qty{0,5}{nm}.\par\smallskip
\editionEtape{D21 · Combien de mesures faut-il ?}
\textbf{Indice 1.} Les quatre nombres affichés sont $n$, la moyenne, l'écart-type $s$ et $u=\dfrac{s}{\sqrt{n}}$. Puisque $\sqrt{100}=10$, comparer $u$ pour 5 et pour
500 mesures.\par\smallskip
\textbf{Indice 2.} De $u=\dfrac{s}{\sqrt{n}}$, on tire $n=\left(\dfrac{s}{u}\right)^2$. Une burette mal réglée décale toutes les valeurs simulées de la même quantité :
regarder ce que devient la moyenne, pas seulement $u$.\par\smallskip
\end{blocEdition}
